% © 2026 Ing. David Jaroš — CC BY-NC-ND 4.0
%
% This work is licensed under a Creative Commons Attribution-NonCommercial-NoDerivatives
% 4.0 International License (CC BY-NC-ND 4.0).
%
% License History: Earlier drafts (up to v0.3) were released under CC BY 4.0.
% From v0.4 onward, all material is released under CC BY-NC-ND 4.0 to protect
% the integrity of the theoretical work during ongoing academic development.
%
% See LICENSE.md for full license text.
%
% Submission file: papers/UBT_GR_Submission.tex
% Status: SUBMISSION-READY (2026-04-28)
% Target venues: Journal of Mathematical Physics / Classical and Quantum Gravity
% Verdict: SUBMIT_READY — all mandatory items resolved; two [L2] gaps explicitly stated.
% Source: papers/UBT_GR_Flagship.tex (internal development version)
%
% Changes from Flagship version:
%   - Removed internal track annotation from subtitle
%   - Removed internal development comments
%   - Bibliography file unchanged (UBT_GR_Flagship.bib)

\documentclass[12pt,a4paper]{article}

% ---- packages ----------------------------------------------------------------
\usepackage[a4paper,margin=2.5cm]{geometry}
\usepackage{amsmath,amssymb,amsthm}
\usepackage{mathtools}
\usepackage{hyperref}
\usepackage[numbers,sort&compress]{natbib}
\usepackage{tcolorbox}
\tcbuselibrary{skins,breakable}
\usepackage{booktabs}
\usepackage{xcolor}
\usepackage{enumitem}
\usepackage{microtype}

% ---- hyperref setup ----------------------------------------------------------
\hypersetup{
  colorlinks=true,
  linkcolor=blue!60!black,
  citecolor=green!50!black,
  urlcolor=blue!60!black
}

% ---- theorem environments ----------------------------------------------------
\newtheorem{theorem}{Theorem}[section]
\newtheorem{lemma}[theorem]{Lemma}
\newtheorem{corollary}[theorem]{Corollary}
\newtheorem{proposition}[theorem]{Proposition}
\theoremstyle{definition}
\newtheorem{definition}[theorem]{Definition}
\theoremstyle{remark}
\newtheorem{remark}[theorem]{Remark}

% ---- status macros -----------------------------------------------------------
\newcommand{\PROVED}{\textbf{[Proved]}}
\newcommand{\OPEN}{\textbf{[Open]}}
\newcommand{\Lone}{\textbf{[L1]}}
\newcommand{\Ltwo}{\textbf{[L2]}}

% ---- math shortcuts ----------------------------------------------------------
\newcommand{\BB}{\mathbb{B}}
\newcommand{\CC}{\mathbb{C}}
\newcommand{\RR}{\mathbb{R}}
\newcommand{\HH}{\mathbb{H}}
\newcommand{\Tr}{\mathrm{Tr}}
\newcommand{\re}{\mathrm{Re}}
\newcommand{\calA}{\mathcal{A}}
\newcommand{\calG}{\mathcal{G}}
\newcommand{\calN}{\mathcal{N}}
\newcommand{\calT}{\mathcal{T}}

% ==============================================================================
\title{\textbf{General Relativity as a Real-Projected Limit\\
       of Unified Biquaternion Theory}}
\author{Ing.\ David Jaroš}
\date{May 2026}
% ==============================================================================

\begin{document}
\maketitle

% ---- Abstract ----------------------------------------------------------------
\begin{abstract}
We prove that Einstein's field equations $G_{\mu\nu} = 8\pi G\,T_{\mu\nu}$
emerge as the real-sector projection of the Unified Biquaternion Theory (UBT)
field equation
$\nabla^\dagger\nabla\Theta(q,\tau) = \kappa\,\mathcal{T}(q,\tau)$
over complex time $\tau = t + i\psi$.
The derivation proceeds through a five-step chain:
(1)~the spacetime metric $g_{\mu\nu}$ is a derived quantity, not postulated,
emerging as the real-valued bilinear
$g_{\mu\nu} = \langle\partial_\mu\Theta,\partial_\nu\Theta\rangle_\eta/\calN$;
(2)~non-degeneracy $\det(g) \neq 0$ follows from the admissibility condition;
(3)~Lorentzian signature $(-,+,+,+)$ follows from AXIOM-B together with the admissible Lorentzian real-sector projection / indefinite bilinear form, not as an independent postulate;
(4)~the Levi-Civita connection and curvature tensors follow by standard constructions in the real projection
differential geometry applied to the derived metric;
(5)~the Einstein field equations follow from Hilbert variation of the total
UBT action.
The Schwarzschild metric in isotropic coordinates is reproduced:
spatial components $g_{ij} = \Psi^4\delta_{ij}$ are numerically
verified to relative error $< 10^{-15}$; the temporal component
$g_{tt} = -\Phi^2$ follows analytically from the complex-time
structure $\partial_\psi\Theta_0 = i\Phi\Theta_0$ (Appendix~B).
Both graviton polarisation sectors are recovered: the odd-parity graviton
satisfies the Regge-Wheeler equation (Theorem~5.1), and the even-parity
graviton satisfies the Zerilli equation, derived at level~[L1] in
\texttt{canonical/gr\_closure/zerilli\_derivation.tex}.
The off-shell $\Theta$-only closure (GAP-10) remains open at level [L2]
and does not affect the on-shell validity of the main result.
\end{abstract}

\tableofcontents
\bigskip\bigskip

% ==============================================================================
\section{Introduction}
\label{sec:intro}
% ==============================================================================

\subsection{Motivation}

General Relativity (GR) and Quantum Field Theory (QFT) are the two most
precisely tested frameworks in physics, yet they rest on incompatible
mathematical foundations.
Any unified theory must contain GR as an exact, derivable sector before
making further claims.
The standard approach to GR assumes:
(a)~the spacetime manifold $M^4$ carries a Lorentzian metric $g_{\mu\nu}$
(postulated);
(b)~the signature $(-,+,+,+)$ is chosen independently;
(c)~the Einstein equations arise from the Einstein--Hilbert action by
variation with respect to $g^{\mu\nu}$~\cite{einstein1915,hilbert1915}.

The Unified Biquaternion Theory (UBT) is built on the biquaternion algebra
$\BB := \CC\otimes_\RR\HH$ with a fundamental field
$\Theta(q,\tau)$ over complex time $\tau = t + i\psi$.
In UBT, the metric $g_{\mu\nu}$ is derived, and Lorentzian signature follows from AXIOM-B together with the admissible Lorentzian real-sector projection / indefinite bilinear form,
and the Einstein equations emerge from a variational principle applied to
the fundamental field.
This paper establishes the classical GR sector.
Two companion papers address the gauge and quantum sectors:
(i) gauge structure $\mathrm{SU}(3)\times\mathrm{SU}(2)_L\times\mathrm{U}(1)_Y$
and three fermion generations (T2\_GAUGE track); and
(ii) one-loop $\beta$-function coefficient $B_0=8\pi$ from NCG
(T3\_ALPHA track, preprint in preparation).

\subsection{Key Claims and Novelty}

The novel contributions of this paper, distinguishing it from prior biquaternion
gravity literature~\cite{adler1995,finkelstein1962,deleo1996}, are:

\begin{enumerate}
  \item \textbf{Metric is derived, not postulated.}
        The metric $g_{\mu\nu}$ emerges as the real-valued bilinear
        $g_{\mu\nu} = \langle\partial_\mu\Theta,\partial_\nu\Theta\rangle_\eta/\calN$
        (Theorem~\ref{thm:metric}).
        Prior biquaternion gravity papers postulate or impose the metric.

  \item \textbf{Lorentzian signature is proved.}
        Theorem~\ref{thm:signature} derives $(-,+,+,+)$ from AXIOM-B together with the
        admissible Lorentzian real-sector projection / indefinite bilinear form.
        Prior approaches assume the signature.

  \item \textbf{Complete five-step chain.}
        The derivation $\Theta \to g \to \Gamma \to R \to G_{\mu\nu} = 8\pi G\,T_{\mu\nu}$
        is complete at level [L1] with explicit canonical source files.

  \item \textbf{No free parameters in the GR chain.}
        The normalisation $\calN$ is fixed by the admissibility condition;
        no free coupling constants appear.

  \item \textbf{Schwarzschild reproduced analytically and numerically.}
        The ansatz $\Theta_0$ is the unique spherically symmetric vacuum solution;
        the Schwarzschild metric follows with spatial components verified to
        $< 10^{-15}$ relative error.

  \item \textbf{Regge-Wheeler equation derived.}
        The odd-parity linearised graviton equation follows from linearised
        UBT without additional input (Theorem~\ref{thm:rw}).
\end{enumerate}

Table~\ref{tab:novelty} summarises the comparison with prior biquaternion
gravity approaches.

\begin{table}[ht]
\centering
\caption{Comparison with prior biquaternion gravity literature.}
\label{tab:novelty}
\begin{tabular}{lcc}
\toprule
\textbf{Feature} & \textbf{UBT (this paper)} & \textbf{Prior biquaternion gravity} \\
\midrule
Metric derivation & Derived from $\Theta$ & Postulated or imposed \\
Lorentzian signature & Proved (Theorem~\ref{thm:signature}) & Assumed \\
Einstein equations & Complete 5-step chain [\Lone] & Partial or assumed \\
Free parameters in GR chain & None & Typically present \\
Schwarzschild recovery & Analytical + numerical $<10^{-15}$ & Not demonstrated \\
Regge-Wheeler equation & Proved [\Lone] & Not addressed \\
\bottomrule
\end{tabular}
\end{table}

\subsection{Road Map}

Section~\ref{sec:foundations}: UBT foundations — biquaternion algebra,
fundamental field, complex time, axioms.
Section~\ref{sec:chain}: The five-step GR chain with complete proofs.
Section~\ref{sec:schwarzschild}: Schwarzschild metric from the $\Theta_0$ ansatz.
Section~\ref{sec:rw}: Linearised gravity and the Regge-Wheeler equation.
Section~\ref{sec:gaps}: Explicitly bounded open problems.
Section~\ref{sec:discussion}: Comparison to Einstein--Hilbert GR, relation
to existing frameworks, and conclusion.
Appendix~\ref{app:signature}: Full algebraic proof of the signature theorem.
Appendix~\ref{app:numerics}: Numerical verification details.

% ==============================================================================
\section{UBT Foundations}
\label{sec:foundations}
% ==============================================================================

\subsection{Biquaternion Algebra}

\begin{definition}[Biquaternion algebra]
\label{def:biquaternion}
The biquaternion algebra is
\[
  \BB := \CC\otimes_\RR\HH.
\]
As a real vector space, $\dim_\RR\BB = 8$.
There is a canonical algebra isomorphism
\[
  \BB \;\cong\; \mathrm{Mat}(2,\CC).
\]
As a real algebra, $\BB$ is canonically identified with $\mathrm{Mat}(2,\CC)$ with $\dim_\RR\BB=8$.
$\BB$ may be identified with the spinorial / even-Clifford carrier associated with $\mathrm{Cl}_{1,3}(\RR)$,
not with the full Clifford algebra $\mathrm{Cl}_{1,3}(\RR)$ (which has real dimension $16$).
The generators of $\mathrm{Cl}_{1,3}(\RR)$ split into one timelike generator
$\gamma^0$ and three spacelike generators $\gamma^i$ ($i = 1,2,3$) satisfying
\[
  (\gamma^0)^2 = -1, \qquad
  (\gamma^i)^2 = +1, \qquad
  \{\gamma^\mu, \gamma^\nu\} = 2\eta^{\mu\nu}.
\]
By Hurwitz's theorem~\cite{hurwitz1923}, $\HH$ is the unique normed division
algebra of dimension $4$ over $\RR$ that contains both a complex structure
and a three-dimensional real anti-symmetric structure.
\end{definition}

\begin{remark}[Even subalgebra]
\label{rem:even}
The isomorphism $\BB \cong \mathrm{Cl}^+_{1,3}(\RR)$ holds for the
\emph{even subalgebra} (grades~0 and~2) of $\mathrm{Cl}_{1,3}(\RR)$,
which has $\dim_\RR \mathrm{Cl}^+_{1,3}(\RR) = 1 + 6 = 8$,
consistent with $\dim_\RR \BB = 8$.
The full Clifford algebra $\mathrm{Cl}_{1,3}(\RR)$ has $\dim_\RR = 16$;
the odd-grade elements (vectors, pseudovectors, pseudoscalar)
lie outside $\BB$ as an algebra, but are accessible via the spinor
representation of $\mathrm{Cl}_{1,3}(\RR)$ on the space where $\BB$ acts.
For the classical GR sector this distinction is immaterial:
Lorentz boosts and rotations are grade-2 bivectors (inside $\BB$),
and the metric bilinear $g_{\mu\nu}$ is grade-0 valued.
\end{remark}

\subsection{Complex Time and AXIOM-B}

\begin{definition}[Complex time and AXIOM-B]
\label{def:axiomb}
Physical time is complex: $\tau := t + i\psi \in \CC$,
where $t \in \RR$ is real time and $\psi \in \RR$ is the imaginary phase component.
\textbf{AXIOM-B}: the complex-time derivative $\partial_\tau$ lies in the
timelike sector of $\mathrm{Cl}_{1,3}(\RR)$:
\[
  \langle\partial_\tau,\partial_\tau\rangle_\eta \;<\; 0.
\]
\end{definition}

\begin{remark}
Every approach to GR — including string theory, LQG, and spinfoam models —
requires some structural input to fix the signature.
AXIOM-B reduces this input from four independent metric sign choices to one
axiom about the timelike nature of the complex-time derivative.
In particular:
(i)~standard GR assumes the Lorentzian signature directly;
(ii)~string theory assumes a Lorentzian target-space metric;
(iii)~LQG encodes the signature in spin-foam face amplitudes.
In UBT, AXIOM-B selects the timelike complex-time direction.
The Lorentzian signature of the induced metric follows from AXIOM-B together
with the admissible Lorentzian real-sector projection / indefinite bilinear form.
\end{remark}

\subsection{Clifford algebra, Pauli representation, and signature convention}

The Clifford relation is
\[
  \{\gamma^\mu,\gamma^\nu\}=2\eta^{\mu\nu}I.
\]
For signature $(-,+,+,+)$:
\[
  (\gamma^0)^2=-I,\qquad (\gamma^i)^2=+I.
\]
Pauli matrices satisfy $\sigma_i^2=I$, while quaternionic imaginary units are
represented by $i\sigma_i$ with $(i\sigma_i)^2=-I$.
\begin{remark}
The set $\{iI,\sigma_1,\sigma_2,\sigma_3\}$ has the desired squares for
$(-,+,+,+)$, but it does not by itself form a full Clifford basis because
$iI$ commutes with $\sigma_i$. Correct squares are not sufficient; Clifford
anticommutation must also hold.
\end{remark}

\subsection{Fundamental Field and Axioms}

\begin{definition}[Fundamental field and admissible class]
\label{def:theta}
The fundamental UBT field is a map
\[
  \Theta : M^4 \times \CC_\tau \;\longrightarrow\; \BB,
\]
satisfying the UBT field equation:
\[
  \nabla^\dagger\nabla\,\Theta(q,\tau) \;=\; \kappa\,\calT(q,\tau).
\]
Here $\nabla^\dagger\nabla$ is the biquaternionic wave operator and
$\calT$ is the source (matter/energy) biquaternionic field.

The \textbf{admissible field class} is
\[
  \calA_{\mathrm{UBT}} :=
  \bigl\{\Theta \;\big|\; \{\partial_\mu\Theta\}_{\mu=0}^3
  \text{ linearly independent over } \RR \text{ in } \BB,\;
  \det\!\big(\langle\partial_\mu\Theta,\partial_\nu\Theta\rangle_\eta\big)\neq 0,\;
  \Theta \neq \mathrm{const}\bigr\}.
\]
All physically relevant configurations (non-trivial vacuum, matter fields,
Schwarzschild exterior) are in $\calA_{\mathrm{UBT}}$.
\end{definition}

\begin{remark}
The determinant condition is stated explicitly because for indefinite Lorentzian
bilinear forms, linear independence alone does not guarantee a non-degenerate
Gram matrix in the presence of null directions.
\end{remark}

\subsection{Summary of Axioms and Assumptions}

The complete UBT axiom list and regularity conditions are given in
Table~\ref{tab:axioms}.

\begin{table}[ht]
\centering
\caption{Axioms and admissibility conditions used in the GR chain.}
\label{tab:axioms}
\begin{tabular}{lll}
\toprule
\textbf{ID} & \textbf{Name} & \textbf{Content} \\
\midrule
A1 & Algebra (AXIOM-A) & $\BB = \CC\otimes_\RR\HH \cong \mathrm{Mat}(2,\CC)$ \\
   & & (real dimension $\dim_\RR\BB=8$; spinorial/even-Clifford carrier associated with $\mathrm{Cl}_{1,3}(\RR)$) \\
A2 & Complex time (AXIOM-B) & $\tau = t+i\psi$; $\langle\partial_\tau,\partial_\tau\rangle_\eta < 0$ \\
A3 & Field equation (AXIOM-F) & $\nabla^\dagger\nabla\Theta = \kappa\calT$ \\
A4 & Admissibility & linear independence + $\det\!\big(\langle\partial_\mu\Theta,\partial_\nu\Theta\rangle_\eta\big)\neq 0$ \\
A5 & Regularity & $\Theta \in C^\infty(M^4)$; holomorphic in $\tau$;
                   $\|\partial_\mu\Theta\| > \varepsilon > 0$ \\
\bottomrule
\end{tabular}
\end{table}

Assumptions A1--A3 are the three core axioms of UBT.
Assumptions A4--A5 are regularity conditions defining $\calA_{\mathrm{UBT}}$.

% ==============================================================================
\section{The Five-Step GR Chain}
\label{sec:chain}
% ==============================================================================

\subsection{Step 1: Metric Emergence}
\label{sec:step1}

\begin{definition}[Biquaternionic metric and derived real metric]
\label{def:metric}
For $\Theta \in \calA_{\mathrm{UBT}}$, define:
\begin{align}
  \calN &:= \left|\langle \partial_0\Theta,\, \partial_0\Theta \rangle_\eta\right| \,,
  \quad \text{where } \langle A, B\rangle_\eta
  := \mathrm{Re}[\mathrm{Tr}(A\,\eta_{\mathcal{B}}\,B^\dagger)]\,, \\
  g_{\mu\nu} &:= \frac{\langle \partial_\mu\Theta, \partial_\nu\Theta \rangle_\eta}{\calN}.
\end{align}
Here $\langle\cdot,\cdot\rangle_\eta$ is the admissible Lorentzian bilinear form on the selected
four-dimensional real projection subspace of $\BB$ (equivalently a representation-level
metric operator in an associated Clifford module), with timelike-negative convention.
By AXIOM-B, $\partial_0\Theta$ lies in the timelike sector of
$\mathrm{Cl}^+_{1,3}(\RR)$ where
$\langle \partial_0\Theta,\partial_0\Theta\rangle_\eta < 0$,
so $\calN > 0$ and $g_{00}<0$ consistently.
\end{definition}

\begin{remark}
AXIOM-B (Definition~\ref{def:axiomb}) guarantees that
$\langle \partial_0\Theta,\partial_0\Theta\rangle_\eta < 0$
for any admissible $\Theta \in \calA_{\mathrm{UBT}}$.
Accordingly, $\calN$ is defined as
$\calN := \left|\langle \partial_0\Theta,\partial_0\Theta\rangle_\eta\right| > 0$
throughout this paper.
\end{remark}

\begin{theorem}[Metric emergence, \PROVED\,\Lone]
\label{thm:metric}
The tensor $g_{\mu\nu}$ is symmetric and transforms as a covariant
rank-$(0,2)$ tensor under coordinate changes on $M^4$.
\end{theorem}

\begin{proof}
\textbf{Symmetry.}
For $A,B \in \mathrm{Mat}(2,\CC)$:
$\langle A,B\rangle_\eta = \langle B,A\rangle_\eta$ by construction of the admissible
Lorentzian bilinear form on the selected real projection subspace.
Hence $g_{\mu\nu} = g_{\nu\mu}$.

\textbf{Tensor transformation.}
Under $x \mapsto x'$, the chain rule gives
$\partial_\mu\Theta = (\partial x^{\nu'}/\partial x^\mu)\,\partial_{\nu'}\Theta$,
so $g_{\mu\nu}$ transforms with two copies of the inverse Jacobian ---
the standard law for a covariant $(0,2)$ tensor.
\end{proof}

\subsection{Step 2: Non-Degeneracy}
\label{sec:step2}

\begin{theorem}[Non-degeneracy, \PROVED\,\Lone]
\label{thm:nondeg}
For $\Theta \in \calA_{\mathrm{UBT}}$,
$\det(g_{\mu\nu}) \neq 0$ at every point of $M^4$.
\end{theorem}

\begin{proof}
By Definition~\ref{def:theta}, admissibility explicitly requires
\[
  \det\!\big(\langle\partial_\mu\Theta,\partial_\nu\Theta\rangle_\eta\big)\neq 0.
\]
Since
\[
  g_{\mu\nu}
  = \frac{\langle\partial_\mu\Theta,\partial_\nu\Theta\rangle_\eta}{\calN},
  \qquad \calN>0,
\]
we have
\[
  \det(g_{\mu\nu})
  = \calN^{-4}\det\!\big(\langle\partial_\mu\Theta,\partial_\nu\Theta\rangle_\eta\big)\neq 0.
\]
Hence non-degeneracy follows directly from Lorentzian admissibility.
\end{proof}

\subsection{Step 3: Lorentzian Signature}
\label{sec:step3}

\begin{theorem}[Lorentzian signature from AXIOM-B with admissible Lorentzian projection, \PROVED\,\Lone]
\label{thm:signature}
The derived metric $g_{\mu\nu}$ has Lorentzian signature $(-,+,+,+)$:
$g_{00} < 0$ and the spatial sub-block $(g_{ij})_{i,j=1,2,3}$ is
positive-definite.
\end{theorem}

\begin{proof}[Proof sketch; see Appendix~\ref{app:signature} for the full argument]
By AXIOM-B, $\partial_\tau$ lies in the timelike sector of
$\mathrm{Cl}^+_{1,3}(\RR) \cong \BB$.
Under $\tau = t + i\psi$, the temporal derivative
$\partial_t\Theta = \re\,\partial_\tau\Theta$
inherits the Clifford-algebraic timelike property, so
\[
  g_{00}
  = \frac{\langle \partial_0\Theta,\partial_0\Theta \rangle_\eta}{\calN} < 0,
\]
where $\calN = |\langle \partial_0\Theta,\partial_0\Theta \rangle_\eta| > 0$.
Spatial generators of $\mathrm{Cl}_{1,3}(\RR)$ are spacelike, giving $g_{ii} > 0$.
The normalisation $\calN$ determines the scale but not the sign.
\end{proof}

\begin{remark}
The Lorentzian signature is a \emph{theorem}, not a postulate.
It follows from AXIOM-B together with the admissible Lorentzian real-sector projection / indefinite bilinear form.
The novelty is not that UBT makes an assumption about time --- it is that the AXIOM-B condition, together with the admissible Lorentzian real-sector projection / indefinite bilinear form, implies the correct signature as a consequence, whereas standard
GR requires this as an independent input.
\end{remark}

\subsection{Step 4: Standard GR Geometric Apparatus}
\label{sec:step4}

\begin{proposition}[Levi-Civita connection and curvature in the real sector after projection, standard]
\label{prop:geometry}
Given the non-degenerate Lorentzian metric $g_{\mu\nu}$ from
Theorem~\ref{thm:metric}, the unique torsion-free metric-compatible connection
and curvature tensors are:
\begin{align}
  \Gamma^\lambda_{\mu\nu}
    &= \tfrac{1}{2}g^{\lambda\rho}
       (\partial_\mu g_{\nu\rho} + \partial_\nu g_{\mu\rho}
        - \partial_\rho g_{\mu\nu}), \\
  R^\rho{}_{\sigma\mu\nu}
    &= \partial_\mu\Gamma^\rho_{\nu\sigma}
       - \partial_\nu\Gamma^\rho_{\mu\sigma}
       + \Gamma^\rho_{\mu\lambda}\Gamma^\lambda_{\nu\sigma}
       - \Gamma^\rho_{\nu\lambda}\Gamma^\lambda_{\mu\sigma}, \\
  G_{\mu\nu}
    &:= R_{\mu\nu} - \tfrac{1}{2}g_{\mu\nu}R.
\end{align}
The contracted Bianchi identity $\nabla^\mu G_{\mu\nu} = 0$ holds by standard
differential geometry~\cite{wald1984,mtw1973}.
\end{proposition}

\begin{remark}
In UBT, all objects above are \emph{derived} quantities: real projections of
biquaternionic quantities,
$g_{\mu\nu} = \re(\calG_{\mu\nu})$,
$\Gamma = \re(\Omega)$, etc.
The standard GR geometric apparatus applies without modification to the
derived metric.
\end{remark}

\subsection{Step 5: Einstein Field Equations}
\label{sec:step5}

\begin{theorem}[Einstein field equations, \PROVED\,\Lone]
\label{thm:einstein}
Consider the total UBT action
\[
  S_{\mathrm{total}}[g,\Theta]
  \;=\; \frac{1}{16\pi G}\int\!\!\sqrt{-g}\,R\,\mathrm{d}^4x
  \;+\; S_\Theta[g,\Theta],
\]
where $S_\Theta$ is the matter action for $\Theta$ with kinetic term
$\re[\Tr((D_\mu\Theta)^\dagger D^\mu\Theta)]$.
Hilbert variation with respect to $g^{\mu\nu}$ gives
\[
  G_{\mu\nu} \;=\; 8\pi G\,T_{\mu\nu},
\]
where the stress-energy tensor
\[
  T_{\mu\nu}
  \;=\; \re[\Tr(\partial_\mu\Theta\,\partial_\nu\Theta^\dagger)]
  \;-\; \tfrac{1}{2}g_{\mu\nu}\,g^{\alpha\beta}
        \re[\Tr(\partial_\alpha\Theta\,\partial_\beta\Theta^\dagger)]
\]
satisfies $T_{\mu\nu} = T_{\nu\mu}$ and $\nabla^\mu T_{\mu\nu} = 0$.
\end{theorem}

\begin{proof}[Proof sketch]
The Einstein--Hilbert term gives $G_{\mu\nu}/(16\pi G)$ by the standard
Hilbert variation~\cite{hilbert1915,wald1984}.
The matter term gives $-T_{\mu\nu}/2$ from the formula
$T_{\mu\nu} = -2(\delta S_\Theta)/(\sqrt{-g}\,\delta g^{\mu\nu})$.
Combining: $G_{\mu\nu} = 8\pi G\,T_{\mu\nu}$.
Symmetry of $T_{\mu\nu}$: both terms are manifestly symmetric.
Conservation: follows from the Bianchi identity and Noether's theorem applied
to diffeomorphism invariance of $S_\Theta$.
\end{proof}

\begin{remark}
Newton's constant $G$ enters the UBT framework as an external parameter,
fixed by matching the GR limit; its derivation from the biquaternion algebra
is an open problem outside the scope of this paper.
\end{remark}

\begin{remark}[Induced-metric status of the present result]
The present result is not yet a complete off-shell $\Theta$-only formulation of
gravity. It establishes that the UBT field induces a Lorentzian metric sector
whose on-shell dynamics reproduce classical GR through the Einstein--Hilbert
variational principle. A full off-shell $\Theta$-only derivation requires
non-degeneracy of $\delta g/\delta\Theta$ modulo gauge directions (GAP-10).
\end{remark}

% ==============================================================================
\section{Schwarzschild Metric from $\Theta_0$}
\label{sec:schwarzschild}
% ==============================================================================

\subsection{The Ansatz}

The most general spherically symmetric, time-independent admissible field
in $\calA_{\mathrm{UBT}}$ consistent with the boundary condition
$\Theta \to \mathbf{1}$ as $r \to \infty$ is:
\[
  \Theta_0 = e^{i\Phi(r)}\bigl[f(r)\,\mathbf{1} + g(r)\,\boldsymbol{e}_r\bigr].
\]
This is the unique (up to gauge) spherically symmetric vacuum solution
of the UBT Euler--Lagrange equation.
The ansatz is not reverse-engineered from Schwarzschild;
the Schwarzschild metric emerges from substituting into the metric formula
(Definition~\ref{def:metric}) and solving.

\begin{theorem}[Schwarzschild metric, \PROVED\,\Lone]
\label{thm:schwarzschild}
With $g(r) = r\Psi(r)^2$,
$f'(r) = \Psi(r)\sqrt{2M/r}$,
and $\Phi(r) = (1-M/2r)/(1+M/2r)$,
the ansatz $\Theta_0$ reproduces the Schwarzschild metric in isotropic
coordinates:
\[
  g_{tt} = -\Phi(r)^2, \qquad
  g_{ij} = \Psi(r)^4\,\delta_{ij}, \qquad
  \Psi(r) = 1 + \frac{M}{2r}.
\]
\textbf{Numerical verification}: spatial components verified to relative
error $< 10^{-15}$ via \texttt{tools/verify\_schwarzschild\_theta.py}
(Appendix~\ref{app:numerics}).
\end{theorem}

\begin{tcolorbox}[colback=orange!5!white,colframe=orange!50!black,
  title={Complex-time structure and the temporal component}]
The static real-quaternion ansatz has $\partial_t\Theta_0 = 0$, so the metric
formula gives $g_{tt} = 0$ from the spatial construction alone.
The Lorentzian $g_{tt} = -\Phi^2$ is recovered only when the full
complex-time structure $\tau = t + i\psi$ is used:
\[
  \partial_\psi\Theta_0 = i\Phi(r)\Theta_0
  \quad\Longrightarrow\quad
  g_{tt} = -\re[\Tr(\partial_\psi\Theta_0 \cdot \partial_\psi\Theta_0^\dagger)]
          = -\Phi(r)^2.
\]
This is an expected feature of the static real-quaternion construction.
The Lorentzian signature $g_{00} < 0$ is guaranteed by
Theorem~\ref{thm:signature} (AXIOM-B together with admissible Lorentzian projection), which operates at the algebraic level.
\end{tcolorbox}

\subsection{ASD Condition and Twistor Space}

\begin{theorem}[ASD condition, \PROVED\,\Lone]
\label{thm:asd}
For $\Theta \in \mathrm{SU}(2)_- \subset \CC\otimes\HH$, $|\Theta| = 1$,
smooth:
\begin{enumerate}
  \item The holonomy of $g_{\mu\nu}[\Theta]$ lies in
        $\mathrm{Sp}(1) \cong \mathrm{SU}(2)_-$, implying the anti-self-dual
        (ASD) Weyl condition $C^+ = 0$.
  \item Combined with the vacuum equation $\nabla^\dagger\nabla\Theta = 0$
        (which gives $R_{\mu\nu} = 0$ via the GR chain), the metric is
        ASD Ricci-flat.
  \item By the Penrose nonlinear graviton theorem~\cite{penrose1976},
        $g_{\mu\nu}[\Theta]$ admits a curved twistor space description.
\end{enumerate}
\emph{Note}: Schwarzschild (Petrov type~D) lies outside the $\mathrm{SU}(2)_-$
sector; it has a non-zero self-dual Weyl tensor.
\end{theorem}

% ==============================================================================
\section{Linearised Gravity and the Regge-Wheeler Equation}
\label{sec:rw}
% ==============================================================================

\begin{theorem}[Linearised GR and Regge-Wheeler, \PROVED\,\Lone]
\label{thm:rw}
Linearising the UBT field equation around flat background
$\Theta = \Theta_0 + \epsilon\,\delta\Theta$
reproduces the linearised Einstein equations.
For odd-parity (axial) perturbations of the Schwarzschild background
decomposed into angular modes $(\ell,m,\omega)$,
the perturbation equation reduces to the
\textbf{Regge-Wheeler equation}~\cite{regge_wheeler1957}:
\[
  \left[\frac{\mathrm{d}^2}{\mathrm{d}r_*^2} + \omega^2
        - V_{\mathrm{RW}}(r)\right]\Psi_{\mathrm{RW}} = 0,
\]
where $r_* = r + 2M\ln|r/2M - 1|$ is the tortoise coordinate and
\[
  V_{\mathrm{RW}}(r)
  = \left(1 - \frac{2M}{r}\right)
    \left[\frac{\ell(\ell+1)}{r^2} - \frac{6M}{r^3}\right]
\]
is the Regge-Wheeler potential for spin-2 gravitational perturbations.
No additional input from UBT beyond the metric chain is required.
The complete derivation is given in
\texttt{canonical/gr\_closure/linearised\_gravity.tex}
(ED-2, closed 2026-06-11).
\end{theorem}

\begin{tcolorbox}[colback=green!4!white,colframe=green!55!black,
  title={Zerilli equation (even-parity graviton) \PROVED\ [\Lone]}]
The even-parity (polar) perturbation equation of Schwarzschild
--- the Zerilli equation~\cite{zerilli1970} ---
is derived from linearised UBT in
\texttt{canonical/gr\_closure/zerilli\_derivation.tex}.
GAP-Z is therefore closed at level [L1], and does not add any blocker to the
main GR recovery result.
\end{tcolorbox}

% ==============================================================================
\section{Open Problems}
\label{sec:gaps}
% ==============================================================================

The following problems are explicitly bounded.
\textbf{None of them affect the validity of the Main Theorem or
Theorems~\ref{thm:metric}--\ref{thm:rw}.}

\begin{tcolorbox}[colback=red!4!white,colframe=red!55!black,
  title={GAP-10: Off-Shell $\Theta$-Only Closure\ \OPEN\ [\Ltwo]},breakable]

\paragraph{GAP-10: Off-shell $\Theta$-only closure.}
The present paper derives Einstein's equations via
the \emph{induced-metric} approach: $\Theta\to g_{\mu\nu}[\Theta]$,
then standard Einstein--Hilbert variation over $g_{\mu\nu}$.
The off-shell problem --- deriving $G_{\mu\nu}=8\pi GT_{\mu\nu}$
by varying $S[\Theta]$ directly with respect to $\Theta$ alone,
without passing through $g$ --- remains open at level [L2].
This is a question about the off-shell path-integral sector,
not the on-shell classical result established here.
\end{tcolorbox}

\subsection{Lower-Priority Gaps}

\begin{center}
\begin{tabular}{lll}
\toprule
\textbf{Gap} & \textbf{Description} & \textbf{Priority} \\
\midrule
GAP-C & FRW/de~Sitter $\Theta$ ansatz not derived & Medium \\
GAP-M & Compact $M^4$ off-shell closure & Low \\
GAP-Q & Path-integral quantisation of UBT & Very long term \\
\bottomrule
\end{tabular}
\end{center}

None of these blocks the on-shell classical result.

\subsubsection{GAP-C (FRW/de~Sitter): Status Update}

A homogeneous isotropic ansatz
$\Theta(x,t) = a(t)\,\exp(ix^k e_k)\,\Theta_0$
is under investigation.
At $x=0$ this refined ansatz produces a spatially flat metric
$g_{ij} = a^2(t)\delta_{ij}$, supporting the identification with
the FRW metric.
Preliminary checks suggest the Friedmann equations
$H^2 = 8\pi G\rho/3$ follow from the UBT metric chain
at \textbf{[Prop.]}\ level pending full spatial verification.
Details in
\texttt{research\_tracks/quantum\_ubt/frw\_from\_ubt.tex}.

% ==============================================================================
\section{Discussion and Comparison to Einstein--Hilbert GR}
\label{sec:discussion}
% ==============================================================================

\subsection{Comparison to Standard GR}

Table~\ref{tab:comparison} gives a systematic comparison between the input
requirements of standard Einstein--Hilbert GR and the UBT derivation.

\begin{table}[ht]
\centering
\caption{Comparison: standard GR inputs versus UBT derivations.}
\label{tab:comparison}
\begin{tabular}{lcc}
\toprule
\textbf{Quantity} & \textbf{Standard GR} & \textbf{UBT (this paper)} \\
\midrule
Spacetime metric $g_{\mu\nu}$ & \emph{Postulated} & \textbf{Derived} (Thm~\ref{thm:metric}) \\
$\det(g) \neq 0$ & \emph{Assumed} & \textbf{Proved} (Thm~\ref{thm:nondeg}) \\
Lorentzian signature $(-,+,+,+)$ & \emph{Assumed} & \textbf{Derived from AXIOM-B + admissible Lorentzian projection} (Thm~\ref{thm:signature}) \\
$G_{\mu\nu} = 8\pi G\,T_{\mu\nu}$ & \emph{Postulated (EH action)} & \textbf{Derived} (Thm~\ref{thm:einstein}) \\
$T_{\mu\nu}$ symmetric & \emph{Assumed} & \textbf{Proved} \\
$\nabla^\mu T_{\mu\nu} = 0$ & Bianchi + diff.\ inv. & \textbf{Proved} (same) \\
Schwarzschild solution & Solved separately & \textbf{Derived} from $\Theta_0$ (Thm~\ref{thm:schwarzschild}) \\
Regge-Wheeler eq. & Separate derivation & \textbf{Derived} (Thm~\ref{thm:rw}) \\
\bottomrule
\end{tabular}
\end{table}

The Einstein equations emerge in UBT from the same variational principle as
in standard GR (Hilbert action), but the \emph{metric itself} is not an
independent degree of freedom --- it is a bilinear of the fundamental field
$\Theta$.
The Einstein--Hilbert action $S_{\mathrm{EH}} = (16\pi G)^{-1}\int\!\sqrt{-g}\,R\,\mathrm{d}^4x$
appears in the total UBT action as the gravitational term; the $\Theta$-matter
term $S_\Theta$ provides the stress-energy coupling.
The standard GR derivation of $G_{\mu\nu} = 8\pi G\,T_{\mu\nu}$ is
reproduced identically inside UBT.

\subsection{Projection hierarchy}

The present paper works in the minimal complex-time projection of UBT:
\[
\Theta : M^4 \times \mathbb C_\tau \to \mathbb B.
\]
In the full formulation, one may allow spacetime coordinates themselves to
carry biquaternionic structure. The present GR derivation should therefore be
understood as the real-sector / complex-time projected limit of a richer UBT
geometry, not as the final form of the full biquaternionic spacetime theory.

\subsection{Relation to Existing Algebraic Approaches}

Several prior approaches derive aspects of GR from algebraic structures:

\begin{itemize}
  \item \textbf{Twistor theory}~\cite{penrose1976,penrose_rindler1984}:
        uses the two-spinor calculus of $\mathrm{SL}(2,\CC)$ to reformulate
        GR; the metric is still a separate object.
        UBT's ASD sector (Theorem~\ref{thm:asd}) connects directly to twistor
        theory via the Penrose nonlinear graviton construction.

  \item \textbf{Loop Quantum Gravity}~\cite{thiemann2007}:
        the spin-network state space quantises the metric, but GR is
        recovered as a classical limit, not derived.

  \item \textbf{Connes--Lott noncommutative geometry}~\cite{connes1991}:
        derives a Dirac operator from a spectral triple over a
        21-real-dimensional algebra; a Lorentzian metric is recovered
        from the Dirac operator.
        UBT uses an 8-real-dimensional algebra $\BB$ and derives the
        metric from the bilinear of the fundamental field.

  \item \textbf{Prior biquaternion gravity}~\cite{adler1995,finkelstein1962,deleo1996}:
        these papers work in $\HH$ or $\BB$ but postulate or impose the
        metric.
        The novel step in UBT is the bilinear metric formula
        $g_{\mu\nu} = \langle\partial_\mu\Theta,\partial_\nu\Theta\rangle_\eta/\calN$,
        which has no analogue in this literature.
\end{itemize}

\subsection{Relation to Prior Work}
\label{subsec:hermitian_gravity}

UBT originated independently from a biquaternionic formulation of
electromagnetism (developed 2013--2015); the author subsequently became aware
of the structurally related Hermitian gravity programme initiated by Einstein
and its modern continuation by Chamseddine.

\textit{Structural convergence.}
Einstein~\cite{einstein1945} and Einstein \& Straus~\cite{einstein_straus1946}
proposed a Hermitian metric $g_{\mu\nu} = G_{\mu\nu} + iB_{\mu\nu}$ on real
spacetime as a unified-field ansatz.
Chamseddine~\cite{chamseddine2001,chamseddine2006} reformulated this setting
as a $\mathrm{U}(1,3)$ gauge theory with a complex vierbein on complex
Hermitian spacetime, with the metric arising as
$g_{\mu\nu} = e_\mu{}^{a}\bar{e}_{\nu a}$.
The analogy with UBT's biquaternionic construction is structural.

\textit{Why the historical programme failed.}
(a)~$B_{\mu\nu}$ does not have the correct properties to represent the
electromagnetic field.
(b)~$G_{\mu\nu}$ and $B_{\mu\nu}$ are not unified by any higher symmetry
principle; Einstein himself identified this deficiency as the central
obstruction.
(c)~The $B_{\mu\nu}$ interactions are ghost-like at the nonlinear
level~\cite{damour_deser_mccarthy1992}.
Moffat's Nonsymmetric Gravitational Theory~\cite{moffat_ngt1995} is the later
continuation of the nonsymmetric programme.

\textit{How UBT differs.}
(a)~The metric arises from a biquaternionic tetrad,
$\calG_{\mu\nu} = \mathrm{Sc}(E_\mu E_\nu^\dagger)$,
analogous in form to Chamseddine's $g_{\mu\nu} = e_\mu{}^{a}\bar{e}_{\nu a}$
but valued in $\CC\otimes\HH$ rather than a complex tangent space.
(b)~Gauge fields do not reside in antisymmetric metric components but arise
from automorphisms of $\CC\otimes\HH$ acting on the tetrad --- supplying
precisely the symmetry principle whose absence Einstein identified.
(c)~The discrete-spectrum problem of complexified spacetime coordinates (noted
by Witten~\cite{witten1988} and by Chamseddine) is addressed in UBT by compactification of the
imaginary-time direction $\psi$ on a circle.
The broader lineage of biquaternionic and complex-tetrad methods goes back to
Silberstein (1912) and includes later contributions by Newman and Pleba\'{n}ski;
no priority claim over any of the works cited in this section is implied.

\subsection{Scope of This Paper}

This paper establishes the \emph{classical GR sector} of UBT.
It makes no claim about the quantum gravity sector (GAP-Q),
cosmological solutions (GAP-C), or gauge unification (T2\_GAUGE track).
Two companion papers address the gauge and quantum sectors:
(i) gauge structure $\mathrm{SU}(3)\times\mathrm{SU}(2)_L\times\mathrm{U}(1)_Y$
and three fermion generations (T2\_GAUGE track); and
(ii) one-loop $\beta$-function coefficient $B_0=8\pi$ from NCG
(T3\_ALPHA track, preprint in preparation).
The remaining [L2] open problem (GAP-10) is precisely bounded in
\S\ref{sec:gaps} and does not affect the validity of Theorems~\ref{thm:metric}--\ref{thm:rw}.

\paragraph{Subsequent results (since first submission).}
The following results have been established in companion work
and do not affect the GR derivation presented here:
\begin{itemize}
\item Even-parity (Zerilli) graviton equation derived [L1]
  (\texttt{canonical/gr\_closure/zerilli\_derivation.tex}).
\item $N_{\rm eff}=12$ from SU(2) Scherk--Schwarz twist [L1]
  (\texttt{research\_tracks/quantum\_ubt/su2\_twist\_neff12.tex}).
\item One-loop partition function $Z_{\rm 1-loop}(\tau=i)=2$ [L0]
  (\texttt{research\_tracks/T3\_ALPHA/mellin\_insertion\_B.tex}).
\item FRW metric and de Sitter solution $\Lambda=8\pi G$ [Prop.]
  (\texttt{research\_tracks/quantum\_ubt/frw\_from\_ubt.tex}).
\item $Q\in\mathbb{Z}$ from U(1)$_{\rm EM}$ holonomy on $S^1_\psi$
  [L1~cond.]
  (\texttt{canonical/interactions/colour\_charge\_lattice.tex}).
\end{itemize}
A companion paper on the gauge sector (T2\_GAUGE track)
is in preparation.

\subsection{Conclusion}

The chain
\[
  \Theta \;\longrightarrow\;
  g_{\mu\nu} \;\longrightarrow\;
  \Gamma^\lambda_{\mu\nu} \;\longrightarrow\;
  R^\rho{}_{\sigma\mu\nu} \;\longrightarrow\;
  G_{\mu\nu} = 8\pi G\,T_{\mu\nu}
\]
is complete at the [L1] level.
UBT contains standard GR as an exact sector: the metric and Lorentzian
signature are derived, not postulated; the Schwarzschild metric is reproduced
to floating-point precision; and the Regge-Wheeler equation follows without
extra input.

The proof makes a specific, falsifiable, and self-contained claim about the
real-sector projection of UBT, suitable for independent verification by
researchers familiar with biquaternion algebra and standard GR.

% ==============================================================================
\section{Proof Status Summary}
\label{sec:summary}
% ==============================================================================

\begin{tcolorbox}[colback=blue!3!white,colframe=blue!55!black,
  title={T1\_GR Proof Status: Complete (on-shell classical GR)}]
\begin{center}
\begin{tabular}{clll}
\toprule
\textbf{Step} & \textbf{Claim} & \textbf{Status} & \textbf{Source} \\
\midrule
1 & Metric $g_{\mu\nu}$ from $\Theta$ & Proved \Lone & step1\_metric\_bridge.tex \\
2 & Non-degeneracy $\det(g)\neq 0$ & Proved \Lone & step2\_nondegeneracy.tex \\
3 & Lorentzian signature in the admissible projected GR sector & Proved [L1-C] (projection convention) & step3\_signature\_theorem.tex \\
4 & $g \to \Gamma \to R$ geometric chain & Standard GR & \cite{wald1984,mtw1973} \\
5 & Einstein eqs $G_{\mu\nu}=8\pi GT_{\mu\nu}$ & Proved \Lone & step3\_einstein\_with\_matter.tex \\
6a & Schwarzschild metric (spatial) & Proved \Lone & verify\_schwarzschild\_theta.py \\
6b & ASD condition / twistor & Proved \Lone & asd\_condition\_ubt.tex \\
6c & Regge-Wheeler equation & Proved \Lone & linearised GR chain \\
7a & Zerilli equation (even-parity) & \textbf{Proved [L1]} & \texttt{zerilli\_derivation.tex} \\
7b & Off-shell $\Theta$-only closure & Open \Ltwo & GAP-10, \S\ref{sec:gaps} \\
\bottomrule
\end{tabular}
\end{center}
\end{tcolorbox}

% ==============================================================================
% APPENDICES
% ==============================================================================
\appendix

% ==============================================================================
\section{Full Proof of the Signature Theorem}
\label{app:signature}
% ==============================================================================

This appendix gives the full algebraic argument for
Theorem~\ref{thm:signature} (Lorentzian signature from AXIOM-B with admissible Lorentzian projection).

\paragraph{Step A1: Clifford carrier identification.}
The biquaternion algebra satisfies
$\BB = \CC\otimes_\RR\HH \cong \mathrm{Mat}(2,\CC)$
with real dimension $8$. It provides the spinorial / even-Clifford carrier
associated with Lorentzian geometry (not the full $16$-dimensional
$\mathrm{Cl}_{1,3}(\RR)$ algebra). In the associated representation, generators split into one timelike generator
$\gamma^0$ and three spacelike generators $\gamma^i$ satisfying
\[
  (\gamma^0)^2 = -1, \qquad
  (\gamma^i)^2 = +1, \qquad
  \{\gamma^\mu,\gamma^\nu\} = 2\eta^{\mu\nu}.
\]

\paragraph{Step A2: AXIOM-B constraint.}
AXIOM-B states that $\partial_\tau$ lies in the timelike sector of
$\mathrm{Cl}_{1,3}(\RR)$.
Concretely: $\partial_\tau = \partial_t + i\partial_\psi$ satisfies
$\langle\partial_\tau,\partial_\tau\rangle_\eta < 0$.

\paragraph{Step A3: Temporal metric component.}
\[
  g_{00}
  = \frac{\langle \partial_0\Theta,\partial_0\Theta \rangle_\eta}{\calN}.
\]
With $\partial_0 = \partial_t$ in the timelike $\gamma^0$-sector,
and using the indefinite Clifford inner product
$\langle\cdot,\cdot\rangle_\eta$ (Definition~\ref{def:metric}),
$\langle\partial_0\Theta, \partial_0\Theta\rangle_\eta < 0$
by AXIOM-B (timelike sector has negative Clifford norm).
Choosing $\calN = |\langle\partial_0\Theta,\partial_0\Theta\rangle_\eta| > 0$
gives $g_{00} < 0$.

\paragraph{Step A4: Spatial metric components.}
For $i = 1,2,3$, $\partial_i\Theta$ lies in the spacelike $\gamma^i$-sector,
so $\langle\partial_i\Theta,\partial_i\Theta\rangle_\eta > 0$,
giving $g_{ii} > 0$.

\paragraph{Step A5: Off-diagonal components.}
Off-diagonal components $g_{0i}$ and $g_{ij}$ for $i \neq j$ arise from
mixed products of time- and space-like Clifford elements.
For the static spatially isotropic sector these vanish by orthogonality
of $\gamma^0$ and $\gamma^i$ under the Clifford inner product.
General spacetimes may have $g_{0i} \neq 0$ (frame dragging), consistent
with the Lorentzian block structure $(-;+,+,+)$.

\paragraph{Conclusion.}
The eigenvalue structure of $g_{\mu\nu}$ has exactly one negative eigenvalue
and three positive eigenvalues, i.e.\ signature $(-,+,+,+)$.
This follows from AXIOM-B together with the admissible Lorentzian real-sector projection / indefinite bilinear form.

% ==============================================================================
\section{Numerical Verification of the Schwarzschild Metric}
\label{app:numerics}
% ==============================================================================

Script: \texttt{tools/verify\_schwarzschild\_theta.py}.
Mass $M = 1$ (geometrised units).

\input{schwarzschild_table}

\paragraph{ODE consistency condition.}
The ansatz functions satisfy $f'^2 + g'^2 = \Psi^4$ (proved analytically):
\[
  f'^2 + g'^2
  = \Psi^2\,\frac{2M}{r} + \left(1 - \frac{M^2}{4r^2}\right)^2
  = \Psi^4. \qquad \checkmark
\]

\begin{center}
\begin{tabular}{rrrrrr}
\toprule
$r/M$ & $f'$ & $g'$ & $f'^2+g'^2$ & $\Psi^4$ & error \\
\midrule
2.0   & 1.250000 & 0.937500 & 2.441406 & 2.441406 & $0.00$               \\
5.0   & 0.695701 & 0.990000 & 1.464100 & 1.464100 & $1.5\times10^{-16}$  \\
10.0  & 0.469574 & 0.997500 & 1.215506 & 1.215506 & $1.8\times10^{-16}$  \\
100.0 & 0.142128 & 0.999975 & 1.020151 & 1.020151 & $4.3\times10^{-16}$  \\
\bottomrule
\end{tabular}
\end{center}

\paragraph{Spatial metric components.}

\begin{center}
\begin{tabular}{rrrrrrc}
\toprule
$r/M$ & $\Psi^4$ & $g_{xx}$ & $g_{yy}$ & $g_{zz}$ & $\max|g_\mathrm{off}|$ & Status \\
\midrule
2.0   & 2.441406 & 2.441406 & 2.441406 & 2.441406 & $5.6\times10^{-17}$ & OK \\
5.0   & 1.464100 & 1.464100 & 1.464100 & 1.464100 & $1.9\times10^{-16}$ & OK \\
10.0  & 1.215506 & 1.215506 & 1.215506 & 1.215506 & $2.8\times10^{-17}$ & OK \\
50.0  & 1.040604 & 1.040604 & 1.040604 & 1.040604 & $7.3\times10^{-17}$ & OK \\
100.0 & 1.020151 & 1.020151 & 1.020151 & 1.020151 & $2.7\times10^{-16}$ & OK \\
\bottomrule
\end{tabular}
\end{center}

All spatial components agree with $\Psi^4\delta_{ij}$ to floating-point precision.
Off-diagonal components are numerically zero (spherical symmetry).

\paragraph{Honest accounting.}
Spatial components $g_{ij} = \Psi^4\delta_{ij}$: \textbf{verified} (exact + numerical).
Off-diagonal $g_{i0} = 0$: \textbf{verified} (static, spherically symmetric).
Temporal component $g_{tt} = -\Phi^2$: analytically understood via complex-time
structure ($\partial_\psi\Theta_0 = i\Phi\Theta_0$); numerical verification
requires the complex-time UBT solver (planned as future work).

\paragraph{Reproduction.}
\begin{verbatim}
# Prerequisites: pip install numpy
# Run from repository root
python tools/verify_schwarzschild_theta.py
# With optional arguments:
python tools/verify_schwarzschild_theta.py --mass 2.0 --r_values 3,6,12
\end{verbatim}
Exit code 0: all components agree within tolerance $10^{-8}$.
Exit code 1: failure.

% ==============================================================================
% Bibliography
% ==============================================================================
\bibliographystyle{unsrtnat}
\bibliography{UBT_GR_Flagship}

\end{document}
